% --- Archon physics formalization source begin ---
% archon:physics
% archon:covers PhyXMiniProblems/problem_phyx_mini_0660.lean
% archon:source-report reports/phyx_mini/problem_phyx_mini_0660.source.json

\chapter{Physics problem phyx\_mini\_0660}
\label{ch:PhyXMiniProblems_problem_phyx_mini_0660}

\paragraph{Problem source.}
\#\# Physical scenario

Alpha decay occurs when an alpha particle tunnels through the Coulomb barrier. Figure shows a simple one-dimensional model of the potential-energy well of an alpha particle in a nucleus with \$A \textbackslash{}approx 235\$. The \$15\$ \$fm\$ width of this one-dimensional potential-energy well is the diameter of the nucleus. Further, to keep the model simple, the Coulomb barrier has been modeled as a \$20\$-\$fm\$-wide, \$30\$-\$MeV\$-high rectangular potential-energy barrier.

\#\# Question

What is the tunneling probability?

\#\# Answer choices

- A: A: \$6.2 \textbackslash{}times 10\textasciicircum{}\{-39\}\$

- B: B: \$6.4 \textbackslash{}times 10\textasciicircum{}\{-39\}\$

- C: C: \$6.6 \textbackslash{}times 10\textasciicircum{}\{-39\}\$

- D: D: \$6.8 \textbackslash{}times 10\textasciicircum{}\{-39\}\$

\#\# Figure caption (auxiliary; use the image as primary evidence)

The image is a graph depicting a potential energy function labeled "U (MeV)" on the y-axis and "x" on the horizontal axis. 

Key features include:

1. **Potential Energy Levels**:
   - Two rectangular shapes represent potential energy barriers, both reaching up to 30 MeV and having a depth of -60 MeV.

2. **Dimensions**: 
   - The widths of the outer rectangles are 20 femtometers (fm) each.
   - The central lower barrier's width is 15 femtometers (fm).

3. **Energy Indicator**:
   - There is a horizontal line labeled "E = 5.0 MeV" that crosses the rectangles, indicating the energy level within the potential.

4. **Measurements**:
   - The 30 MeV line is above the x-axis, while the -60 MeV line is below it, indicating the depth and height of the potential wells.

This setup is often used to illustrate potential wells or barriers in quantum mechanics, such as those found in nuclear physics.

\#\# Dataset metadata

Quantum Phenomena; Physical Model Grounding Reasoning, Numerical Reasoning

\paragraph{Recorded answer/context.}
C: C: \$6.6 \textbackslash{}times 10\textasciicircum{}\{-39\}\$

\paragraph{Figure/image path.}
/root/proposal\_for\_physic/science-mango/phyx\_mini\_run/phyx\_data/test\_image/660.png

\paragraph{Formalization target.}
create a compiling Lean file with sorry bodies at `PhyXMiniProblems/problem\_phyx\_mini\_0660.lean`. The Lean declarations must preserve the physical quantities, dimensions or dimensional roles, figure labels, governing-law hypotheses, and final relation expressed by this problem.
Use Mathlib/Physlib names found through LeanExplore where available. If a physics API is missing, introduce faithful local abstractions rather than scalar placeholder aliases.

\begin{theorem}[Physics formalization target]
\label{thm:physics:phyx_mini_0660:target}
\lean{PhyXMiniProblems.ProblemPhyXMini0660.problem_phyx_mini_0660}
\uses{def:physics:phyx-mini-0660:phyxminiproblems-problemphyxmini0660-energyinmegaelectronvolts, def:physics:phyx-mini-0660:phyxminiproblems-problemphyxmini0660-alphadecaybarriersetup, def:physics:phyx-mini-0660:phyxminiproblems-problemphyxmini0660-matchesalphadecayscenario, def:physics:phyx-mini-0660:phyxminiproblems-problemphyxmini0660-matchesproblemprosereadouts, def:physics:phyx-mini-0660:phyxminiproblems-problemphyxmini0660-matchessuppliednuclearpotentialfigure, def:physics:phyx-mini-0660:phyxminiproblems-problemphyxmini0660-usesstandardalphareferencedata, def:physics:phyx-mini-0660:phyxminiproblems-problemphyxmini0660-matchesphyslibrectangularbarriers, def:physics:phyx-mini-0660:phyxminiproblems-problemphyxmini0660-hasphysicalalphatunnelingparameters, def:physics:phyx-mini-0660:phyxminiproblems-problemphyxmini0660-satisfiesalphabarrierdecayrelation, def:physics:phyx-mini-0660:phyxminiproblems-problemphyxmini0660-satisfiesopaquealphatransmissionlaw, def:physics:phyx-mini-0660:phyxminiproblems-problemphyxmini0660-answerchoice, def:physics:phyx-mini-0660:phyxminiproblems-problemphyxmini0660-roundstodisplayedscientificnotation, def:physics:phyx-mini-0660:phyxminiproblems-problemphyxmini0660-matchesdisplayedtunnelingprobability}
The assigned autoformalize agent should translate this physics problem into Lean declarations in the covered file, with theorem and lemma proof bodies written as `by sorry`.
\end{theorem}
\begin{proof}
This is an autoformalization task, not a proof task. Produce statements that can later be proved by the physics prover without weakening the physical model.
\end{proof}

% --- Archon declaration topology begin ---
\section{Declaration topology}
\begin{definition}[Length Quantity]
  \label{def:physics:phyx-mini-0660:phyxminiproblems-problemphyxmini0660-lengthquantity}
  \lean{PhyXMiniProblems.ProblemPhyXMini0660.LengthQuantity}
  A nonnegative, unit-independent physical length
\end{definition}
\begin{proof}
  This is the declared model definition of Length Quantity.
\end{proof}
\begin{definition}[Mass Quantity]
  \label{def:physics:phyx-mini-0660:phyxminiproblems-problemphyxmini0660-massquantity}
  \lean{PhyXMiniProblems.ProblemPhyXMini0660.MassQuantity}
  A nonnegative, unit-independent physical mass
\end{definition}
\begin{proof}
  This is the declared model definition of Mass Quantity.
\end{proof}
\begin{definition}[Energy Quantity]
  \label{def:physics:phyx-mini-0660:phyxminiproblems-problemphyxmini0660-energyquantity}
  \lean{PhyXMiniProblems.ProblemPhyXMini0660.EnergyQuantity}
  A signed, unit-independent physical energy
\end{definition}
\begin{proof}
  This is the declared model definition of Energy Quantity.
\end{proof}
\begin{definition}[action Dimension]
  \label{def:physics:phyx-mini-0660:phyxminiproblems-problemphyxmini0660-actiondimension}
  \lean{PhyXMiniProblems.ProblemPhyXMini0660.actionDimension}
  The physical dimension mass * length² / time of action
\end{definition}
\begin{proof}
  This is the declared model definition of action Dimension.
\end{proof}
\begin{definition}[Action Quantity]
  \label{def:physics:phyx-mini-0660:phyxminiproblems-problemphyxmini0660-actionquantity}
  \lean{PhyXMiniProblems.ProblemPhyXMini0660.ActionQuantity}
  \uses{def:physics:phyx-mini-0660:phyxminiproblems-problemphyxmini0660-actiondimension}
  A nonnegative, unit-independent physical action
\end{definition}
\begin{proof}
  This is the declared model definition of Action Quantity.
\end{proof}
\begin{definition}[Inverse Length Quantity]
  \label{def:physics:phyx-mini-0660:phyxminiproblems-problemphyxmini0660-inverselengthquantity}
  \lean{PhyXMiniProblems.ProblemPhyXMini0660.InverseLengthQuantity}
  A nonnegative inverse length used for under-barrier attenuation
\end{definition}
\begin{proof}
  This is the declared model definition of Inverse Length Quantity.
\end{proof}
\begin{definition}[length Readout]
  \label{def:physics:phyx-mini-0660:phyxminiproblems-problemphyxmini0660-lengthreadout}
  \lean{PhyXMiniProblems.ProblemPhyXMini0660.lengthReadout}
  \uses{def:physics:phyx-mini-0660:phyxminiproblems-problemphyxmini0660-lengthquantity}
  Read a physical length in a selected named unit
\end{definition}
\begin{proof}
  This is the declared model definition of length Readout.
\end{proof}
\begin{definition}[length In Meters]
  \label{def:physics:phyx-mini-0660:phyxminiproblems-problemphyxmini0660-lengthinmeters}
  \lean{PhyXMiniProblems.ProblemPhyXMini0660.lengthInMeters}
  \uses{def:physics:phyx-mini-0660:phyxminiproblems-problemphyxmini0660-lengthquantity, def:physics:phyx-mini-0660:phyxminiproblems-problemphyxmini0660-lengthreadout}
  Read a physical length in metres
\end{definition}
\begin{proof}
  This is the declared model definition of length In Meters.
\end{proof}
\begin{definition}[length In Femtometers]
  \label{def:physics:phyx-mini-0660:phyxminiproblems-problemphyxmini0660-lengthinfemtometers}
  \lean{PhyXMiniProblems.ProblemPhyXMini0660.lengthInFemtometers}
  \uses{def:physics:phyx-mini-0660:phyxminiproblems-problemphyxmini0660-lengthquantity, def:physics:phyx-mini-0660:phyxminiproblems-problemphyxmini0660-lengthreadout}
  Read a physical length in femtometres
\end{definition}
\begin{proof}
  This is the declared model definition of length In Femtometers.
\end{proof}
\begin{definition}[mass In Kilograms]
  \label{def:physics:phyx-mini-0660:phyxminiproblems-problemphyxmini0660-massinkilograms}
  \lean{PhyXMiniProblems.ProblemPhyXMini0660.massInKilograms}
  \uses{def:physics:phyx-mini-0660:phyxminiproblems-problemphyxmini0660-massquantity}
  Read a physical mass in kilograms
\end{definition}
\begin{proof}
  This is the declared model definition of mass In Kilograms.
\end{proof}
\begin{definition}[energy In Joules]
  \label{def:physics:phyx-mini-0660:phyxminiproblems-problemphyxmini0660-energyinjoules}
  \lean{PhyXMiniProblems.ProblemPhyXMini0660.energyInJoules}
  \uses{def:physics:phyx-mini-0660:phyxminiproblems-problemphyxmini0660-energyquantity}
  Read an energy in coherent SI joules
\end{definition}
\begin{proof}
  This is the declared model definition of energy In Joules.
\end{proof}
\begin{definition}[energy In Electron Volts]
  \label{def:physics:phyx-mini-0660:phyxminiproblems-problemphyxmini0660-energyinelectronvolts}
  \lean{PhyXMiniProblems.ProblemPhyXMini0660.energyInElectronVolts}
  \uses{def:physics:phyx-mini-0660:phyxminiproblems-problemphyxmini0660-energyquantity}
  Read an energy in electron-volts using Physlib's calibrated unit
\end{definition}
\begin{proof}
  This is the declared model definition of energy In Electron Volts.
\end{proof}
\begin{definition}[energy In Mega Electron Volts]
  \label{def:physics:phyx-mini-0660:phyxminiproblems-problemphyxmini0660-energyinmegaelectronvolts}
  \lean{PhyXMiniProblems.ProblemPhyXMini0660.energyInMegaElectronVolts}
  \uses{def:physics:phyx-mini-0660:phyxminiproblems-problemphyxmini0660-energyquantity, def:physics:phyx-mini-0660:phyxminiproblems-problemphyxmini0660-energyinelectronvolts}
  Read an energy in mega-electron-volts
\end{definition}
\begin{proof}
  This is the declared model definition of energy In Mega Electron Volts.
\end{proof}
\begin{definition}[action In Joule Seconds]
  \label{def:physics:phyx-mini-0660:phyxminiproblems-problemphyxmini0660-actioninjouleseconds}
  \lean{PhyXMiniProblems.ProblemPhyXMini0660.actionInJouleSeconds}
  \uses{def:physics:phyx-mini-0660:phyxminiproblems-problemphyxmini0660-actionquantity}
  Read an action in coherent SI joule-seconds
\end{definition}
\begin{proof}
  This is the declared model definition of action In Joule Seconds.
\end{proof}
\begin{definition}[inverse Length In Inverse Meters]
  \label{def:physics:phyx-mini-0660:phyxminiproblems-problemphyxmini0660-inverselengthininversemeters}
  \lean{PhyXMiniProblems.ProblemPhyXMini0660.inverseLengthInInverseMeters}
  \uses{def:physics:phyx-mini-0660:phyxminiproblems-problemphyxmini0660-inverselengthquantity}
  Read an attenuation wave number in inverse metres
\end{definition}
\begin{proof}
  This is the declared model definition of inverse Length In Inverse Meters.
\end{proof}
\begin{definition}[Particle Species]
  \label{def:physics:phyx-mini-0660:phyxminiproblems-problemphyxmini0660-particlespecies}
  \lean{PhyXMiniProblems.ProblemPhyXMini0660.ParticleSpecies}
  Particle species represented by the energy line
\end{definition}
\begin{proof}
  This is the declared model definition of Particle Species.
\end{proof}
\begin{definition}[Barrier Origin]
  \label{def:physics:phyx-mini-0660:phyxminiproblems-problemphyxmini0660-barrierorigin}
  \lean{PhyXMiniProblems.ProblemPhyXMini0660.BarrierOrigin}
  Physical origin assigned to the modeled barriers
\end{definition}
\begin{proof}
  This is the declared model definition of Barrier Origin.
\end{proof}
\begin{definition}[Barrier Side]
  \label{def:physics:phyx-mini-0660:phyxminiproblems-problemphyxmini0660-barrierside}
  \lean{PhyXMiniProblems.ProblemPhyXMini0660.BarrierSide}
  The two barriers flanking the nuclear well in the one-dimensional graph
\end{definition}
\begin{proof}
  This is the declared model definition of Barrier Side.
\end{proof}
\begin{definition}[Potential Region]
  \label{def:physics:phyx-mini-0660:phyxminiproblems-problemphyxmini0660-potentialregion}
  \lean{PhyXMiniProblems.ProblemPhyXMini0660.PotentialRegion}
  The three finite-width regions drawn in the graph
\end{definition}
\begin{proof}
  This is the declared model definition of Potential Region.
\end{proof}
\begin{definition}[Plot Axis]
  \label{def:physics:phyx-mini-0660:phyxminiproblems-problemphyxmini0660-plotaxis}
  \lean{PhyXMiniProblems.ProblemPhyXMini0660.PlotAxis}
  The two axes of the supplied potential-energy plot
\end{definition}
\begin{proof}
  This is the declared model definition of Plot Axis.
\end{proof}
\begin{definition}[Plot Axis Quantity]
  \label{def:physics:phyx-mini-0660:phyxminiproblems-problemphyxmini0660-plotaxisquantity}
  \lean{PhyXMiniProblems.ProblemPhyXMini0660.PlotAxisQuantity}
  Physical role assigned to each plot axis
\end{definition}
\begin{proof}
  This is the declared model definition of Plot Axis Quantity.
\end{proof}
\begin{definition}[Potential Profile Shape]
  \label{def:physics:phyx-mini-0660:phyxminiproblems-problemphyxmini0660-potentialprofileshape}
  \lean{PhyXMiniProblems.ProblemPhyXMini0660.PotentialProfileShape}
  Qualitative shape of the plotted potential
\end{definition}
\begin{proof}
  This is the declared model definition of Potential Profile Shape.
\end{proof}
\begin{definition}[figure Region]
  \label{def:physics:phyx-mini-0660:phyxminiproblems-problemphyxmini0660-barrierside-figureregion}
  \lean{PhyXMiniProblems.ProblemPhyXMini0660.BarrierSide.figureRegion}
  \uses{def:physics:phyx-mini-0660:phyxminiproblems-problemphyxmini0660-barrierside, def:physics:phyx-mini-0660:phyxminiproblems-problemphyxmini0660-potentialregion}
  Region of the primary figure corresponding to a barrier side
\end{definition}
\begin{proof}
  This is the declared model definition of figure Region.
\end{proof}
\begin{definition}[Nuclear Potential Figure]
  \label{def:physics:phyx-mini-0660:phyxminiproblems-problemphyxmini0660-nuclearpotentialfigure}
  \lean{PhyXMiniProblems.ProblemPhyXMini0660.NuclearPotentialFigure}
  \uses{def:physics:phyx-mini-0660:phyxminiproblems-problemphyxmini0660-lengthquantity, def:physics:phyx-mini-0660:phyxminiproblems-problemphyxmini0660-energyquantity, def:physics:phyx-mini-0660:phyxminiproblems-problemphyxmini0660-potentialregion, def:physics:phyx-mini-0660:phyxminiproblems-problemphyxmini0660-plotaxis, def:physics:phyx-mini-0660:phyxminiproblems-problemphyxmini0660-plotaxisquantity, def:physics:phyx-mini-0660:phyxminiproblems-problemphyxmini0660-potentialprofileshape}
  Typed geometric and energy information carried by the supplied graph. No tunneling probability or answer choice is stored in the figure
\end{definition}
\begin{proof}
  This is the declared model definition of Nuclear Potential Figure.
\end{proof}
\begin{definition}[Alpha Decay Barrier Setup]
  \label{def:physics:phyx-mini-0660:phyxminiproblems-problemphyxmini0660-alphadecaybarriersetup}
  \lean{PhyXMiniProblems.ProblemPhyXMini0660.AlphaDecayBarrierSetup}
  \uses{def:physics:phyx-mini-0660:phyxminiproblems-problemphyxmini0660-lengthquantity, def:physics:phyx-mini-0660:phyxminiproblems-problemphyxmini0660-massquantity, def:physics:phyx-mini-0660:phyxminiproblems-problemphyxmini0660-energyquantity, def:physics:phyx-mini-0660:phyxminiproblems-problemphyxmini0660-actionquantity, def:physics:phyx-mini-0660:phyxminiproblems-problemphyxmini0660-inverselengthquantity, def:physics:phyx-mini-0660:phyxminiproblems-problemphyxmini0660-particlespecies, def:physics:phyx-mini-0660:phyxminiproblems-problemphyxmini0660-barrierorigin, def:physics:phyx-mini-0660:phyxminiproblems-problemphyxmini0660-barrierside, def:physics:phyx-mini-0660:phyxminiproblems-problemphyxmini0660-nuclearpotentialfigure}
  Independent physical quantities in the alpha-decay model. In particular, the attenuation wave number and tunneling probability are fields constrained only by the general laws below, not definitions of the recorded answer
\end{definition}
\begin{proof}
  This is the declared model definition of Alpha Decay Barrier Setup.
\end{proof}
\begin{definition}[Matches Alpha Decay Scenario]
  \label{def:physics:phyx-mini-0660:phyxminiproblems-problemphyxmini0660-matchesalphadecayscenario}
  \lean{PhyXMiniProblems.ProblemPhyXMini0660.MatchesAlphaDecayScenario}
  \uses{def:physics:phyx-mini-0660:phyxminiproblems-problemphyxmini0660-alphadecaybarriersetup}
  Qualitative physical model stated in the problem
\end{definition}
\begin{proof}
  This is the declared model definition of Matches Alpha Decay Scenario.
\end{proof}
\begin{definition}[Matches Problem Prose Readouts]
  \label{def:physics:phyx-mini-0660:phyxminiproblems-problemphyxmini0660-matchesproblemprosereadouts}
  \lean{PhyXMiniProblems.ProblemPhyXMini0660.MatchesProblemProseReadouts}
  \uses{def:physics:phyx-mini-0660:phyxminiproblems-problemphyxmini0660-lengthinfemtometers, def:physics:phyx-mini-0660:phyxminiproblems-problemphyxmini0660-energyinmegaelectronvolts, def:physics:phyx-mini-0660:phyxminiproblems-problemphyxmini0660-alphadecaybarriersetup}
  Numerical information stated in the prose. The interval formalizes the source's approximate mass-number description A ≈ 235. This premise has no tunneling-probability field
\end{definition}
\begin{proof}
  This is the declared model definition of Matches Problem Prose Readouts.
\end{proof}
\begin{definition}[Matches Supplied Nuclear Potential Figure]
  \label{def:physics:phyx-mini-0660:phyxminiproblems-problemphyxmini0660-matchessuppliednuclearpotentialfigure}
  \lean{PhyXMiniProblems.ProblemPhyXMini0660.MatchesSuppliedNuclearPotentialFigure}
  \uses{def:physics:phyx-mini-0660:phyxminiproblems-problemphyxmini0660-lengthinfemtometers, def:physics:phyx-mini-0660:phyxminiproblems-problemphyxmini0660-energyinmegaelectronvolts, def:physics:phyx-mini-0660:phyxminiproblems-problemphyxmini0660-barrierside, def:physics:phyx-mini-0660:phyxminiproblems-problemphyxmini0660-barrierside-figureregion, def:physics:phyx-mini-0660:phyxminiproblems-problemphyxmini0660-alphadecaybarriersetup}
  Labels, geometry, and energy levels read directly from image 660.png. The graph shows two equal rectangular barriers, a central well whose width is the nuclear diameter, the exterior zero level, and the alpha energy line
\end{definition}
\begin{proof}
  This is the declared model definition of Matches Supplied Nuclear Potential Figure.
\end{proof}
\begin{definition}[Uses Standard Alpha Reference Data]
  \label{def:physics:phyx-mini-0660:phyxminiproblems-problemphyxmini0660-usesstandardalphareferencedata}
  \lean{PhyXMiniProblems.ProblemPhyXMini0660.UsesStandardAlphaReferenceData}
  \uses{def:physics:phyx-mini-0660:phyxminiproblems-problemphyxmini0660-massinkilograms, def:physics:phyx-mini-0660:phyxminiproblems-problemphyxmini0660-actioninjouleseconds, def:physics:phyx-mini-0660:phyxminiproblems-problemphyxmini0660-alphadecaybarriersetup}
  Standard alpha-particle and reduced-Planck reference data needed for the textbook estimate. These calibrations do not contain a transmission result
\end{definition}
\begin{proof}
  This is the declared model definition of Uses Standard Alpha Reference Data.
\end{proof}
\begin{definition}[Matches Physlib Rectangular Barriers]
  \label{def:physics:phyx-mini-0660:phyxminiproblems-problemphyxmini0660-matchesphyslibrectangularbarriers}
  \lean{PhyXMiniProblems.ProblemPhyXMini0660.MatchesPhyslibRectangularBarriers}
  \uses{def:physics:phyx-mini-0660:phyxminiproblems-problemphyxmini0660-lengthinmeters, def:physics:phyx-mini-0660:phyxminiproblems-problemphyxmini0660-massinkilograms, def:physics:phyx-mini-0660:phyxminiproblems-problemphyxmini0660-energyinjoules, def:physics:phyx-mini-0660:phyxminiproblems-problemphyxmini0660-alphadecaybarriersetup}
  Calibrate each Physlib scalar rectangular barrier against the dimensionful alpha mass, width, and barrier energy. Absolute face coordinates are not fixed because only the widths are relevant to the tunneling calculation
\end{definition}
\begin{proof}
  This is the declared model definition of Matches Physlib Rectangular Barriers.
\end{proof}
\begin{definition}[Has Physical Alpha Tunneling Parameters]
  \label{def:physics:phyx-mini-0660:phyxminiproblems-problemphyxmini0660-hasphysicalalphatunnelingparameters}
  \lean{PhyXMiniProblems.ProblemPhyXMini0660.HasPhysicalAlphaTunnelingParameters}
  \uses{def:physics:phyx-mini-0660:phyxminiproblems-problemphyxmini0660-lengthinmeters, def:physics:phyx-mini-0660:phyxminiproblems-problemphyxmini0660-massinkilograms, def:physics:phyx-mini-0660:phyxminiproblems-problemphyxmini0660-energyinjoules, def:physics:phyx-mini-0660:phyxminiproblems-problemphyxmini0660-actioninjouleseconds, def:physics:phyx-mini-0660:phyxminiproblems-problemphyxmini0660-inverselengthininversemeters, def:physics:phyx-mini-0660:phyxminiproblems-problemphyxmini0660-alphadecaybarriersetup}
  Positivity and branch conditions selecting a physical tunneling setup
\end{definition}
\begin{proof}
  This is the declared model definition of Has Physical Alpha Tunneling Parameters.
\end{proof}
\begin{definition}[Satisfies Alpha Barrier Decay Relation]
  \label{def:physics:phyx-mini-0660:phyxminiproblems-problemphyxmini0660-satisfiesalphabarrierdecayrelation}
  \lean{PhyXMiniProblems.ProblemPhyXMini0660.SatisfiesAlphaBarrierDecayRelation}
  \uses{def:physics:phyx-mini-0660:phyxminiproblems-problemphyxmini0660-massinkilograms, def:physics:phyx-mini-0660:phyxminiproblems-problemphyxmini0660-energyinjoules, def:physics:phyx-mini-0660:phyxminiproblems-problemphyxmini0660-actioninjouleseconds, def:physics:phyx-mini-0660:phyxminiproblems-problemphyxmini0660-inverselengthininversemeters, def:physics:phyx-mini-0660:phyxminiproblems-problemphyxmini0660-alphadecaybarriersetup}
  For E < U, the classically forbidden wave number satisfies κ = sqrt (2 m (U - E)) / ℏ. Both sides are coherent-SI inverse-metre readouts. No numerical transmission probability appears in this law
\end{definition}
\begin{proof}
  This is the declared model definition of Satisfies Alpha Barrier Decay Relation.
\end{proof}
\begin{definition}[Satisfies Opaque Alpha Transmission Law]
  \label{def:physics:phyx-mini-0660:phyxminiproblems-problemphyxmini0660-satisfiesopaquealphatransmissionlaw}
  \lean{PhyXMiniProblems.ProblemPhyXMini0660.SatisfiesOpaqueAlphaTransmissionLaw}
  \uses{def:physics:phyx-mini-0660:phyxminiproblems-problemphyxmini0660-lengthinmeters, def:physics:phyx-mini-0660:phyxminiproblems-problemphyxmini0660-inverselengthininversemeters, def:physics:phyx-mini-0660:phyxminiproblems-problemphyxmini0660-alphadecaybarriersetup}
  The opaque rectangular-barrier estimate T = exp (-2 κ L) for the barrier encountered by the outgoing alpha particle. This is a governing relation among independent quantities, not the requested numerical specialization
\end{definition}
\begin{proof}
  This is the declared model definition of Satisfies Opaque Alpha Transmission Law.
\end{proof}
\begin{definition}[Answer Choice]
  \label{def:physics:phyx-mini-0660:phyxminiproblems-problemphyxmini0660-answerchoice}
  \lean{PhyXMiniProblems.ProblemPhyXMini0660.AnswerChoice}
  Labels of the four dimensionless probability choices
\end{definition}
\begin{proof}
  This is the declared model definition of Answer Choice.
\end{proof}
\begin{definition}[displayed Tunneling Probability]
  \label{def:physics:phyx-mini-0660:phyxminiproblems-problemphyxmini0660-displayedtunnelingprobability}
  \lean{PhyXMiniProblems.ProblemPhyXMini0660.displayedTunnelingProbability}
  \uses{def:physics:phyx-mini-0660:phyxminiproblems-problemphyxmini0660-answerchoice}
  Probability printed beside each answer label, represented exactly
\end{definition}
\begin{proof}
  This is the declared model definition of displayed Tunneling Probability.
\end{proof}
\begin{definition}[recorded Dataset Answer]
  \label{def:physics:phyx-mini-0660:phyxminiproblems-problemphyxmini0660-recordeddatasetanswer}
  \lean{PhyXMiniProblems.ProblemPhyXMini0660.recordedDatasetAnswer}
  \uses{def:physics:phyx-mini-0660:phyxminiproblems-problemphyxmini0660-answerchoice}
  Answer metadata recorded by the source dataset; not a physics premise
\end{definition}
\begin{proof}
  This is the declared model definition of recorded Dataset Answer.
\end{proof}
\begin{definition}[Rounds To Displayed Scientific Notation]
  \label{def:physics:phyx-mini-0660:phyxminiproblems-problemphyxmini0660-roundstodisplayedscientificnotation}
  \lean{PhyXMiniProblems.ProblemPhyXMini0660.RoundsToDisplayedScientificNotation}
  value rounds to the probability reportedValue, where the rounding quantum is 0.1 × 10⁻³⁹ = 10⁻⁴⁰. The half-open interval resolves the upper tie
\end{definition}
\begin{proof}
  This is the declared model definition of Rounds To Displayed Scientific Notation.
\end{proof}
\begin{definition}[Matches Displayed Tunneling Probability]
  \label{def:physics:phyx-mini-0660:phyxminiproblems-problemphyxmini0660-matchesdisplayedtunnelingprobability}
  \lean{PhyXMiniProblems.ProblemPhyXMini0660.MatchesDisplayedTunnelingProbability}
  \uses{def:physics:phyx-mini-0660:phyxminiproblems-problemphyxmini0660-alphadecaybarriersetup, def:physics:phyx-mini-0660:phyxminiproblems-problemphyxmini0660-answerchoice, def:physics:phyx-mini-0660:phyxminiproblems-problemphyxmini0660-displayedtunnelingprobability, def:physics:phyx-mini-0660:phyxminiproblems-problemphyxmini0660-roundstodisplayedscientificnotation}
  A displayed choice is the rounded value of the modeled probability
\end{definition}
\begin{proof}
  This is the declared model definition of Matches Displayed Tunneling Probability.
\end{proof}
% --- Archon declaration topology end ---
% --- Archon physics formalization source end ---
